\clearpage
\chapter*{강연 236의 정오표 및 보충}
\addcontentsline{toc}{chapter}{강연 236의 정오표 및 보충}
\src{303}

\noindent\textbf{8. TDTE VI: 피카르 스킴. 일반 성질.}

\begin{center}
[부르바키 세미나, 제14권, 1961/62, 제236호, 23쪽.]
\end{center}

\noindent\textbf{12쪽, 16행.} --- ``아핀 직선'' 대신
``원점을 뺀 아핀 직선''으로 읽어야 한다.

\noindent\textbf{18행.} --- ``$X[t]$'' 대신
``$X[t,t^{-1}]$''로 읽어야 한다.

\noindent\textbf{16쪽, 명제 3.1, 6행.} --- ``그런데 가정은 …'' 대신
``그런데 가정으로부터 …''로 읽어야 한다.

\noindent\textbf{21쪽, § 4: 피카르 스킴에 대한 유한성 정리.} ---
이 절에서 제기한 종류의 유한성 문제들은 이 강연을 작성한 뒤 거의 완전히
해결되었다. 이제 이 방향에서 알려진 주요 사실들을 제시하자. 명제를
간단히 하기 위하여, 명제들에 나타나는 모든 피카르 준스킴이 존재한다고
암묵적으로 가정한다. 다만 이 명제들을 명백하게 수정하면 어떠한 존재
가정도 둘 필요가 없다. 이하에서 $S$는 뇌터 스킴을, $X,Y$는 $S$ 위에서
고유한 스킴들을 나타낸다.

\medskip
\noindent(i) $f:X\longrightarrow Y$를 전사 사상이라 하자. 그러면
\[
f^*: \operatorname{Pic}_{Y/S}\longrightarrow\operatorname{Pic}_{X/S}
\]
는 유한형 사상이다.

\src{304}
\noindent(ii) $Y$가 $S$ 위에서 사영이고 $S$에 관해 풍부한 가역 가군을
갖는다고 하자. $X$를 이 가군의 임의의 단면의 영점 스킴이라 하고,
$f:X\longrightarrow Y$를 표준 몰입 사상이라 하자. 끝으로 $Y/S$의
섬유들의 기약 성분들이 차원 $\geq3$이라고 가정하자. 그러면
\[
f^*: \operatorname{Pic}_{Y/S}\longrightarrow\operatorname{Pic}_{X/S}
\]
는 유한형이다.

\noindent(iii) $X$가 $S$ 위에서 사영이고 그 모든 기하적 섬유가 차원
$n$인 정역 스킴이라고 하자. $\mathcal O_X(1)$을 $X$ 위의 풍부한 가역
가군으로서 힐베르트 다항식들을 정의할 수 있게 하는 것으로 잡자. $M$을
$\operatorname{Pic}_{X/S}$의 한 부분집합이라 하자. 그러면 $M$이
준콤팩트일 필요충분조건은 $M$의 원소들의 힐베르트 다항식
\[
a_0x^n+a_1x^{n-1}+a_2x^{n-2}+\cdots
\]
에서 계수 $a_1$과 $a_2$가 유계로 남는 것이다. 여기서는 그 단항식의
차수가 음수가 되는 모든 계수를 생략한다. 따라서 $n=1$이면 $a_1$만
관여하고, $n=0$이면 둘 다 관여하지 않는다.

\noindent(iv) 모든 정수 $n\neq0$에 대하여, 군 준스킴
$\operatorname{Pic}_{X/S}$의 $n$제곱 준동형은 유한형 사상이다.

(i)과 (ii)는 또한 (각각의 명제에서 고려한 가정 아래)
$\operatorname{Pic}_{Y/S}$의 부분집합 $M$이 준콤팩트일 필요충분조건이
$\operatorname{Pic}_{X/S}$ 안에서 그 상이 준콤팩트인 것임을 뜻한다.
따라서 $Y$ 위의 가역 가군 $\mathcal L$이 $0$과 $\tau$-동치일
필요충분조건은 $X$로의 그 역상이 그러한 것이다. 다시 말해
$\operatorname{Pic}_{Y/S}^\tau$는
$\operatorname{Pic}_{X/S}^\tau$의 역상이다. 특히 앞의 준스킴이 $S$
위에서 유한형임을 보이려면 뒤의 준스킴에 대하여 이를 증명하면 충분하다.
실제로 $\operatorname{Pic}_{Y/S}\longrightarrow
\operatorname{Pic}_{X/S}$는 유한형이다. 이제 (i), 차우 보조정리와
(iii)을 사용하면 다음을 얻는다.

\medskip
\noindent(v) $\operatorname{Pic}_{X/S}^\tau$는 $S$ 위에서 유한형이다.

일반적으로 유한 사상에 대한 (i)과 (ii)를 함께 사용하면 대부분의 유한성
문제를 $X/S$가 차원 $\leq2$인 기하적으로 정역이고 정규인 섬유들을 갖는
경우로 환원할 수 있다. 흔히 유한일 필요가 없는 전사 사상에 (i)을
적용하고 대수 곡면의 특이점 해소(임의의 표수에서 아브향카르에 의한
결과)를 사용하면, $X/S$가 매끄러워서 기하적 섬유들이 차원 2인
비특이 스킴인 경우로까지 환원된다. 예를 들어 (v)와 비특이 사영 곡면의
네론--세베리 군의 계수(階數)에 상계를 주는 피카르--이구사 부등식을
함께 사용하면, 네론의 유한성 정리의 다음 일반화를 증명할 수 있다.

\medskip
\noindent(vi) $X/S$가 $S$ 위에서 고유하고 그 밖에는 임의적이라고 하자.
그러면 $X/S$의 기하적 섬유 $X_i/k_i$의 네론--세베리 군
\[
\operatorname{Pic}_{X_i/k_i}/\operatorname{Pic}_{X_i/k_i}^\circ
\]

\src{305}
은 유한 생성이며, 그 계수(階數)와 꼬임 부분군의 위수에는 공통 상계가
있다.

비특이 곡면의 경우로 환원하는 같은 방법과 이 경우에 알려진 정리들, 즉
네론의 유한성 정리와 네론--세베리 군 위의 교차 형식이 비퇴화라는 사실은
다음을 함의한다.

\medskip
\noindent(vii) $X$를 대수적으로 닫힌 체 위의 고유 스킴이라 하자.
그러면 $X$ 안에는 다음 성질을 갖는 유한 개의 닫힌 정역 곡선
$C_i$ ($1\leq i\leq r$)가 존재한다. $\operatorname{Pic}_{X/k}$의
모든 부분집합 $M$에 대하여, $M$이 준콤팩트일 필요충분조건은 정수들
\[
\deg_{C_i'}\mathcal L\qquad(\mathcal L\in M)
\]
이 유계로 남는 것이다. 여기서 $C_i'$는 $C_i$의 정규화를 나타낸다.

여기서 $r$은 네론--세베리 군의 계수(階數)로 잡을 수 있다. 이 군이 유한
생성임을 알고 나면, (vii)은 $X$ 안의 곡선 $C$가 네론--세베리 군 위에
정의하는 선형 형식들이 동시에 영이 되는 원소는 네론--세베리 군의 꼬임
원소뿐이라는 사실로 환원된다. $X$가 사영이고 비특이인 경우에는 이 결과와
(v)가 마쓰사카에 의한 것이었다. (vii)을 사용하면 $X$ 위에서 $0$과
$\tau$-동치인 가역 가군들의 다음 특성화를 쉽게 얻는다.

\medskip
\noindent(viii) $X/k$를 한 체 위의 고유 스킴이라 하고, $\mathcal L$을
$X$ 위의 가역 가군이라 하자. 다음 조건들은 서로 동치이다.

\noindent a. $\mathcal L$은 $0$과 $\tau$-동치이다.

\noindent b. 모든 연접 가군 $F$에 대하여
\[
\chi(F\otimes\mathcal L)=\chi(F)
\]
이다. 여기서 $\chi$는 오일러--푸앵카레 특성을 나타낸다.

\noindent b'. (b)에서와 같되 $F=\mathcal O_Y$로 한다. 여기서 $Y$는
$X$ 안의 차원 1인 닫힌 정역 부분스킴이다.

\noindent c. 위와 같은 모든 $Y$에 대하여, $Y'$을 그 정규화 곡선이라
하면
\[
\deg\mathcal L_{Y'}=0
\]
이다.

$X/k$가 사영이고 풍부한 가역 가군 $\mathcal O_X(1)$을 갖는다면, 앞의
조건들은 다음 조건들과도 서로 동치이다.

\noindent d. 모든 정수 $m$에 대하여 $\mathcal L^{\otimes m}(1)$은
풍부하다.

\noindent e. ($X$가 정역인 경우.) 모든 정수의 쌍 $m,n$에 대하여
\[
\chi(\mathcal L^{\otimes m}(n))=\chi(\mathcal O_X(n))
\]

\src{306}
이다. 즉 $F=\mathcal L^{\otimes m}(n)$으로 놓으면 (b)가 성립한다.

마지막 조건이 충분함을 보려면, 이 조건이 $\mathcal L^{\otimes m}$의
힐베르트 다항식들이 모두 같다는 뜻임에 유의한다. 따라서 멈퍼드의 판정법
(iii)에 따라 $\mathcal L^{\otimes m}$들은
$\operatorname{Pic}_{X/k}$의 한 준콤팩트 부분집합 안에 머물며, 이는
(a)이다. 조건 (b), (b'), (c)는 보통 비특이 사영 다양체 위에서 정의하는
수치적 동치 개념을 임의의 고유 스킴 위로 옮긴 변형으로 보아야 한다.
비특이 사영 다양체의 경우 (a)와 (c)의 동치는 물론 잘 알려져 있었다
(마쓰사카).

판정법 (e)는 또한 다음 결과를 함의한다.

\medskip
\noindent(ix) $f:X\longrightarrow S$를 기하적 섬유들이 정역인 사영
평탄 사상이라 하자. 그러면 $\operatorname{Pic}_{X/S}^\tau$는
$\operatorname{Pic}_{X/S}$ 안에서 열리고 닫혀 있다.

핵심 결과 (i), (ii), (iii)의 증명에 관한 몇 가지 논평만 하기로 한다.
결과 (iv)는 다른 결과들과 다소 별개이며, 유한하고 전사이며 라디시엘인
사상, 더 정확히는 프로베니우스 사상에 대하여 (i)만을 사용해서 증명된다.
(i)에는 비평탄 하강의 발상들을 본질적으로 사용한다(특히 TDTE I, 9쪽을
보라). 유한성 결과만을 목표로 하므로 하강 자료의 유효성 판정법이 없는
것은 아무런 문제가 되지 않는다. 한편 멈퍼드는 최근 (iii)의 약간 약한
형태, 즉 힐베르트 다항식의 모든 계수를 사용하는 판정법을 증명했다. 그의
논증은 극히 간단하며, 나카이의 풍부성 판정법의 증명에서 착안한 것이다
(나카이는 비특이 곡면에 대하여 이를 진술했고, 멈퍼드는 임의의 사영
사상으로 일반화했다). 다만 내가 보기에는 이 논증이 $X/S$의 섬유에
약간의 추가 제한, 곧 세르의 조건 $(S_2)$를 둘 때에만 타당하다. 기하적
섬유들이 정규이면 이 조건은 성립한다. 이제 이 제한된 판정법을 (ii)의
증명에 사용한다. 실제로 판정법 (i)에 의해 $Y/S$가 평탄하고 기하적으로
정역이며 정규인 섬유들을 갖는 경우로 환원할 수 있고, 멈퍼드의 판정법을
적용하면 $X/S$가 같은 조건들을 만족하는 경우로 쉽게 환원할 수 있다.
그러면 차원에 대한 가정으로부터 $Y$와 $X$의 기하적 섬유들은 그 닫힌점들
에서 깊이 $\geq2$임을 알 수 있다. 따라서 SGA 1962(강연 XII)에 제시된
형태의 ``동치 판정법''을 적용하여 (ii)의 증명을 끝낼 수 있다. (i)과
(ii)를 얻은 뒤에는 멈퍼드의 판정법에서 섬유들의 정규성에 관한 모든
가정을 제거하고, (iii)에 제시된 더 강한 형태로도 이를 증명하기 어렵지
않다.

\src{307}
끝으로 (i)과 (ii)의 증명은 $S$가 체 $k$의 스펙트럼인 경우 사상
\[
\operatorname{Pic}_{Y/k}\longrightarrow\operatorname{Pic}_{X/k}
\]
이 유한형일 뿐만 아니라 아핀임을 보여 준다.

% FGA Canon R30: scan restoration and explicit editorial apparatus.
\par\smallskip
\begingroup\footnotesize
\noindent\textit{FGA-CANON-DISP-0059 (C125). 인쇄본의 독법은 계수
$a_1$과 $a_2$가 유계로 남아야 한다고 요구하지만 $n<2$인 경우를
명시하지 않는다. 현행 독법은 낮은 차원에서 존재하지 않는 계수에 관한
관례만을 명시한다. 이는 인쇄된 명제가 거짓이라고 판정한 것이 아니다.}\par
\endgroup
